\section*{\texorpdfstring{\S\CurrentExerciseGroup}{Section \CurrentExerciseGroup}}
\phantomsection
\addcontentsline{toc}{section}{\texorpdfstring{Exercises for \S\CurrentExerciseGroup}{Exercises for Section \CurrentExerciseGroup}}

\begin{enumerate}
\def\labelenumi{\arabic{enumi})}
\tightlist
\item
  Show that if \(f\) and \(g\) are two numerical functions \(\geqslant 0\) defined on \(X\) , and \(\mu\) is a positive measure on \(X\) , then
\end{enumerate}

\[
\mu^ {*} (\sup (f, g)) + \mu^ {*} (\inf (f, g)) \leqslant \mu^ {*} (f) + \mu^ {*} (g).
\]

From this, deduce that if A and B are any subsets of X, then

\[
\mu^ {*} (\mathrm{A} \cup \mathrm{B}) + \mu^ {*} (\mathrm{A} \cap \mathrm{B}) \leqslant \mu^ {*} (\mathrm{A}) + \mu^ {*} (\mathrm{B})
\]

(cf.~§4, Exer. 8 d)).

\begin{enumerate}
\def\labelenumi{\arabic{enumi})}
\setcounter{enumi}{1}
\item
  For every \(x \in \mathbf{X}\) , show that for every numerical function \(f \geqslant 0\) defined on \(\mathbf{X}\) , \(\varepsilon_x^*(f) = f(x)\) .
\item
  Give an example, in \(\mathbf{R}\) , of an open set that is not relatively compact, whose outer measure (for Lebesgue measure) is finite.
\item
  Let \(\mathrm{X}\) be a locally compact space, \(\alpha\) a finite numerical function \(\geqslant 0\) on \(\mathrm{X}\) such that, for every compact subset \(\mathrm{K}\) of \(\mathrm{X}\) , \(\sum_{x\in \mathrm{K}}\alpha (x)\) is finite; let \(\mu\) be the positive measure
\end{enumerate}

on X defined by the masses \(\alpha(x)\) (Ch. III, §1, No.~3, Example I).

\begin{enumerate}
\def\labelenumi{\alph{enumi})}
\tightlist
\item
  Show that for every function \(f \geqslant 0\) , lower semi-continuous on \(X\) , one has \(\mu^{*}(f) = \sum_{x \in X} \alpha(x)f(x)\) , with the convention that \(\alpha(x)f(x) = 0\) when \(\alpha(x) = 0\) and \(f(x) = -\infty\) .
\end{enumerate}

\[
f (x) = + \infty
\]

\begin{enumerate}
\def\labelenumi{\alph{enumi})}
\setcounter{enumi}{1}
\tightlist
\item
  Let \(f \geqslant 0\) be any numerical function defined on \(X\) . Show that if \(\mu^{*}(f) < +\infty\) , then \(\mu^{*}(f) = \sum_{x \in X} \alpha(x)f(x)\) , with the same convention as in a). (Cf. Exer. 5.)
\end{enumerate}

¶ 5) Let X be the subset of the plane \(\mathbf{R}^2\) formed by the union of the line \(D = \{0\} \times \mathbf{R}\) and the set of points \((1/n, k/n^2)\) , where \(n\) runs over the set of integers \(>0\) , and \(k\) over the set \(\mathbf{Z}\) of rational integers.

\begin{enumerate}
\def\labelenumi{\alph{enumi})}
\item
  For every point \((0,y)\) of D and every integer n \textgreater{} 0, let \(\mathrm{T}_{n}(y)\) be the set of points \((u,v)\) in X such that \(u \leqslant 1/n\) and \(|v - y| \leqslant u\) . Show that if one takes as fundamental system of neighborhoods of each point \((0,y)\) in D the set of \(\mathrm{T}_{n}(y)\) (n \textgreater{} 0), and as fundamental system of neighborhoods of each of the other points of X the unique set reduced to that point, one defines on X a topology T for which X is a locally compact space that is not paracompact.
\item
  Let \(\alpha\) be the numerical function on \(X\) that is equal to 0 on \(D\) and to \(1/n^3\) at each of the points \((1/n, k/n^2)\) . Show that, for every compact subset \(K\) of \(X\) , \(\sum_{x \in K} \alpha(x)\)
\end{enumerate}

is finite; let \(\mu\) be the positive measure on X defined by the masses \(\alpha(x)\) . Show that \(\mu^{*}(D) = +\infty\) even though \(\alpha(x) = 0\) on D. (If an open set U for \(\mathcal{T}\) contains D, show that there exist an interval \(a \leqslant y \leqslant b\) on D not reduced to a point, a set B dense in this interval (for the usual topology of \(\mathbf{R}\) ), and an integer \(n > 0\) such that, for every \(y \in B\) , one has \(\mathrm{T}_{n}(y) \subset \mathrm{U}\) ; for this, one may make use of Baire's theorem.)

\begin{enumerate}
\def\labelenumi{\arabic{enumi})}
\setcounter{enumi}{5}
\tightlist
\item
  Let \(f\) be a numerical function \(\geqslant 0\) defined on \(X\) .
\end{enumerate}

\begin{enumerate}
\def\labelenumi{\alph{enumi})}
\item
  Show that, for the mapping \(\mu \mapsto \mu^{*}(f)\) of \(\mathcal{M}_{+}(X)\) into \(\overline{\mathbf{R}}\) to be continuous for the vague topology, it is necessary (and sufficient) that \(f\) be continuous with compact support (make use of Exer. 2). For \(\mu \mapsto \mu^{*}(f)\) to be lower semi-continuous for the vague topology, it is necessary (and sufficient, cf.~Prop. 4) that \(f\) be lower semi-continuous.
\item
  Show that, for the mapping \(\mu \mapsto \mu^{*}(f)\) of \(\mathcal{M}_{+}(\mathrm{X})\) into \(\overline{\mathbf{R}}\) to be continuous for the quasi-strong topology (Ch. III, §1, Exer. 8), it is necessary and sufficient that \(f\) be bounded and have compact support (method analogous to that of \(a\) )). From this, deduce that for every function \(f \geqslant 0\) that is zero on the complement of a countable union of compact sets, the mapping \(\mu \mapsto \mu^{*}(f)\) is lower semi-continuous for the quasi-strong topology (make use of Th. 3) (cf.~Exer. 7 b)).
\item
  Show that, for the mapping \(\mu \mapsto \mu^{*}(f)\) of \(\mathcal{M}_{+}(\mathrm{X}) \cap \mathcal{M}^{1}(\mathrm{X})\) into \(\overline{\mathbf{R}}\) to be continuous for the ultrastrong topology (Ch. III, §1, Exer. 15), it is necessary and sufficient that \(f\) be bounded; for every function \(f \geqslant 0\) defined on \(\mathrm{X}\) , \(\mu \mapsto \mu^{*}(f)\) is lower semicontinuous for the ultrastrong topology.
\item
  Show that, for the mapping \(\mu \mapsto \mu^{*}(f)\) of \(\mathcal{M}_{+}(\mathrm{X}) \cap \mathcal{M}^{1}(\mathrm{X})\) into \(\overline{\mathbf{R}}\) to be continuous for the weak topology (Ch. III, §1, Exer. 15), it is necessary and sufficient that \(f\) be continuous on \(\mathrm{X}\) and tend to 0 at the point at infinity; for \(\mu \mapsto \mu^{*}(f)\) to be lower semi-continuous for the weak topology, it is necessary and sufficient that \(f\) be lower semi-continuous.
\end{enumerate}

\begin{enumerate}
\def\labelenumi{\arabic{enumi})}
\setcounter{enumi}{6}
\tightlist
\item
  \begin{enumerate}
  \def\labelenumii{\alph{enumii})}
  \tightlist
  \item
    Let \((\mu_{n})\) be an increasing sequence of measures \(\geqslant0\) on a locally compact space X; assume that the sequence is bounded above in \(\mathcal{M}_{+}(X)\) and denote by \(\mu\) its supremum. Let f be a function \(\geqslant0\) defined on X and zero on the complement of a countable union of compact sets. Show that \(\mu^{*}(f)=\sup\mu_{n}^{*}(f)\) (cf.~Exer. 6 b)).
  \end{enumerate}
\end{enumerate}

\begin{enumerate}
\def\labelenumi{\alph{enumi})}
\setcounter{enumi}{1}
\tightlist
\item
  Let X and \(\mu\) be the locally compact space and the measure defined in Exer. 5. Let \(\alpha_{n}\) be the numerical function equal to \(\alpha\) (in the notation of Exer. 5) for every point \((1/m, k/m^{2})\) such that \(m \leqslant n\) , and equal to 0 at the other points of X; let \(\mu_{n}\) be the measure defined by the masses \(\alpha_{n}(x)\) . Show that \(\mu\) is the supremum of the sequence \((\mu_{n})\) in \(\mathcal{M}_{+}(X)\) and that \(\mu_{n}^{*}(D)=0\) for all n, but \(\mu^{*}(D)=+\infty\) .
\end{enumerate}

\begin{enumerate}
\def\labelenumi{\arabic{enumi})}
\setcounter{enumi}{7}
\tightlist
\item
  \begin{enumerate}
  \def\labelenumii{\alph{enumii})}
  \tightlist
  \item
    Let \((\mu_n)\) be a decreasing sequence of measures \(\geqslant 0\) on a locally compact space \(X\) , and let \(\mu\) be the infimum of the sequence in \(\mathcal{M}_+(X)\) . Show that if \(f\) is a function \(\geqslant 0\) such that \(\mu_n^*(f) < +\infty\) from some index onward, then \(\mu^*(f) = \inf_n \mu_n^*(f)\) (when \(g\) is lower semi-continuous, positive, and satisfies \(\mu^*(g) < +\infty\) , note that there exists a sequence \((h_m)\) of continuous functions \(\geqslant 0\) with compact support, such that \(\sum_{m=1}^{\infty} h_m \leqslant g\) and \(\mu_n^*(g) = \sum_{m=1}^{\infty} \mu_n^*(h_m)\) for every index \(n\) .
  \end{enumerate}
\end{enumerate}

\begin{enumerate}
\def\labelenumi{\alph{enumi})}
\setcounter{enumi}{1}
\tightlist
\item
  On the discrete space X = N, let \(\mu_{n}\) be the measure defined by placing mass +1 at each point \(m \geqslant n\) . Show that the infimum of the decreasing sequence \((\mu_{n})\) in \(\mathcal{M}_{+}(X)\) is 0, but that \(\mu_{n}^{*}(X) = +\infty\) for all n.
\end{enumerate}

\setcounter{exercisegroup}{3}
\edef\CurrentExerciseGroup{\arabic{exercisegroup}}
